
%\needspace{4\baselineskip}
\color{uniblue}{\subsection{\raggedright{Контроллер присоединения (КП)}}}
\color{black}

\begin{enumerate}[label=\arabic{section}.\arabic{subsection}.\arabic{enumi}, labelsep=4pt, leftmargin=0pt, itemindent=57pt, itemsep=0pt, parsep=5pt] % если что можно крутить itemsep

%%%%%%%%%%%%%%%%%%%%%%%%%%%%%%%%%%%%%%%%%%%%%%%%%%%%%%%%%%%%%%%%%%%%%%  ОБЩИЕ СВЕДЕНИЯ %%%%%%%%%%%%%%%%%%%%%%%%%%%%%%%%%%%%%%%%%%%%%%%%%%
\item Общие сведения

\begin{enumerate}[label=\arabic{section}.\arabic{subsection}.\arabic{enumi}.\arabic*, labelsep=4pt, leftmargin=0em, itemindent=65pt, parsep=0pt]

\item
КП предназначен:
\begin{itemize}
\item для оперативного управления коммутационными аппаратами; 
\item для формирования сигнала положения коммутационного аппарата для целей АСУ ТП; 
\item для реализации оперативной блокировки управляемых коммутационных аппаратов.
\end{itemize}

\item 
В устройстве контроллер присоединения представлен функциональным блоком «КП». В функциональном блоке заложена возможность оперативного управления следующими коммутационными аппаратами:
\begin{itemize}
	\item силовой выключатель;
	\item заземляющий нож;
	\item выкатной элемент.
\end{itemize}

\item
При выполнении коммутационной операции В, ЗН, или ВЭ информация об этом процессе будет отображаться общим выходным сигналом «КП: Идет переключение». Функциональная схема логики формирования общего сигнала представлена на рисунке \ref{kp:img_perkl}.

\item
Ввод функционального блока в работу на этапе параметрирования устройства осуществляется программным переключателем SGF1 «Ввод\_функции». Введенное состояние функционального блока соответствует активному состоянию сигнала «КП: Ввод».

\vspace{1mm}
\begin{figure}[H]
\centering
\includegraphics[width=0.6\textwidth,height=0.6\textheight,keepaspectratio]{\fbpath/128. КП ОЛ/_latex/img/2.pdf}
\captionof{figure}{Функциональная схема формирования общего сигнала переключения коммутационного аппарата}\label{kp:img_perkl}
\end{figure}

\item
Функциональный блок может быть выведен из работы оперативно путем активации сигнала «ОВ КП», о чем сигнализирует активное состояние сигнала «КП: Оперативный вывод».

\item
Логика управления режимами работы функционального блока выполнена в соответствии с рисунком \ref{kp:img_gen}. В таблице \ref{kp:tbl_gen} приведены параметры, необходимые для настройки функционального блока. 

\vspace{1mm}
\begin{figure}[H]
\centering
\includegraphics[width=1\textwidth,height=1\textheight,keepaspectratio]{\fbpath/128. КП ОЛ/_latex/img/1.pdf}
\captionof{figure}{Функциональная схема логики управления режимами работы КП}\label{kp:img_gen}
\end{figure}

\end{enumerate}
%%%%%%%%%%%%%%%%%%%%%%%%%%%%%%%%%%%%%%%%%%%%%%%%%%%%%%%%%%%%%%%%%%%%%%  КОНЕЦ ОБЩИЕ СВЕДЕНИЯ %%%%%%%%%%%%%%%%%%%%%%%%%%%%%%%%%%%%%%%%%%%%%%%%%%

%%%%%%%%%%%%%%%%%%%%%%%%%%%%%%%%%%%%%%%%%%%%%%%%%%%%%%%%%%%%%%%%%%%%%%  УПРАВЛЕНИЕ ВЫКЛЮЧАТЕЛЕМ %%%%%%%%%%%%%%%%%%%%%%%%%%%%%%%%%%%%%%%%%%%%%%%%%%

\item Управление выключателем (УВ)

\begin{enumerate}[label=\arabic{section}.\arabic{subsection}.\arabic{enumi}.\arabic*, labelsep=4pt, leftmargin=0em, itemindent=65pt, parsep=0pt]

\item
В функциональном блоке «КП» логика управления выключателем представлена функцией «УВ». Функция формирует команды на отключение и включение выключателя при оперативном управлении.

\item
Местное управление коммутационным аппаратом осуществляется только посредством кнопок \tikz[baseline={(0,-0.00ex)}]{\draw[line width=1.6pt] (0,0) rectangle (1em,1em);\node[anchor=center] at (0.5em,0.5em) {\textbf{I}};}~–~«Включить» и \tikz[baseline={(0,-0.00ex)}]{\draw[line width=1.6pt] (0,0) rectangle (1em,1em);\node[anchor=center] at (0.5em,0.5em) {\textbf{O}};}~–~«Отключить». При этом необходимо перевести устройство в режим «Местное» с помощью заданной функциональной кнопки. В данном случае управление по локальной сети и с помощью выносных кнопок или ключей будет невозможно.
\item
Дистанционное управление может реализоваться через дискретные входы, которые сконфигурированы как «Включить В от ТУ» и «Отключить В от ТУ», а также по локальной сети. В дистанционном режиме управления команды по локальной сети предусмотрены постоянно.
\item
Управление выключателем также осуществляется с помощью пульта управления (ПУ), который предназначен для формирования команд управления выключателем путем автоматической подачи на катушки электромагнитов электрического напряжения. Управление с ПУ может реализоваться через дискретные входы, которые сконфигурированы как «Включить В от ПУ» и «Отключить В от ПУ».

\item
Команда оперативного отключения «УВ: Отключить» формируется в следующих случаях:
\begin{itemize}
	\item при появлении команды на оперативное отключение выключателя («Отключить В~от~ПУ» или «Откл. В~ИЧМ») и выбранным управлением по месту (сигнал «Местное» находится в активном состоянии);
	\item при появлении команды на оперативное отключение выключателя («Отключить В~от~ТУ» или «Откл. В~АСУ») и выбранным дистанционным управлением (сигнал «Дистанционное» находится в активном состоянии).
\end{itemize}
\item
Прохождение команды на отключение блокируется при наличии следующих сигналов:
\begin{itemize}
	\item запрет отключения от функции контроля силового выключателя «КСВ: Блокировка откл»;
	\item отключенное положение силового выключателя «В отключен (б/к)» (с контролем сигнала «В включен~(б/к)»);
	\item превышение времени переключения силового выключателя «УВ: Прев. времени~пер.»;
	\item местное управление силовым выключателем (положение переключателя «М/Д» в приводе выключателя) «Ключ М/Д привода В».
\end{itemize}
\item
В устройстве применяется режим включения выключателя с контролем синхронизма, который контролируется уставкой SGF2 «Контроль\_синхр». Если уставка SGF2 введена в работу, то при наличии разрешающих условий для включения выключателя и активного сигнала разрешения оперативного включения из функции КС «Разреш.опер.вкл.В» формируется команда включения с контролем синхронизма.

\item
Команда оперативного включения «УВ: Включить» без контроля синхронизма формируется в следующих случаях:
\begin{itemize}
	\item при появлении команды на оперативное включение выключателя («Включить В~от~ПУ» или «Вкл. В~ИЧМ») и выбранным управлением по месту (сигнал «Местное» находится в активном состоянии);
	\item при появлении команды на оперативное включение выключателя («Включить В~от~ТУ» или «Вкл. В~АСУ») и выбранным дистанционным управлением (сигнал «Дистанционное» находится в активном состоянии).
\end{itemize}

\item
Прохождение команды на включение блокируется при наличии следующих сигналов:
\begin{itemize}
	\item запрет включения от функции контроля силового выключателя «КСВ: Блокировка вкл»;
	\item включенное положение силового выключателя «В~включен (б/к)» (с контролем сигнала «В~отключен~(б/к)»);
	\item превышение времени переключения силового выключателя «УВ: Прев. времени~пер.»;
	\item местное управление силовым выключателем (положение переключателя «М/Д» в приводе выключателя) «Ключ М/Д привода В»;
	\item появление сигнала аварийного отключения силового выключателя «КСВ: В аварийно отключен» (прохождение сигнала определяется состоянием программного переключателя SGF3 «Бл\_вкл\_от\_авар\_откл»).
\end{itemize}

\item  
При наличии введенной уставки включения выключателя с контролем синхронизма SGF2 и превышение времени регулируемой уставки «T3» формируется выходной сигнал задержки включения для сброса команды.

\item
При наличии сигналов «УВ: Включить» и «УВ: Отключить» длительностью больше, чем время уставки «T1» формируется сигнал превышения допустимого времени переключения -- «УВ: Прев. времени пер.». Активное состояние операции включения или отключения сигнализируется выходным сигналом «УВ: Идет переключение».

\item
Логика функции формирует сигналы состояния выключателя «УВ: Отключено», «УВ: Включено», «УВ: Не определено» и «УВ: Неисправность», при этом сигналы «УВ: Неисправность» и «УВ: Не определено» выдаются с выдержкой времени, задаваемой уставкой Т2 «Тблк».

\item
Ввод функции в работу на этапе параметрирования устройства осуществляется программным переключателем SGF1 «Ввод\_функции». Информация о введенной функции отображается выходным сигналом «УВ: Ввод».

\item
Алгоритм работы выполнен в соответствии с рисунком \ref{kp:img_uv}. 

\vspace{1mm}
\begin{figure}[H]
\centering
\includegraphics[width=1\textwidth,height=1\textheight,keepaspectratio]{\fbpath/128. КП ОЛ/_latex/img/uv.pdf}
\captionof{figure}{Функциональная схема логики УВ}
\label{kp:img_uv}
\end{figure}

\end{enumerate}
%%%%%%%%%%%%%%%%%%%%%%%%%%%%%%%%%%%%%%%%%%%%%%%%%%%%%%%%%%%%%%%%%%%%%% КОНЕЦ УПРАВЛЕНИЕ ВЫКЛЮЧАТЕЛЕМ %%%%%%%%%%%%%%%%%%%%%%%%%%%%%%%%%%%%%%%%%%%%%%%%%%

%%%%%%%%%%%%%%%%%%%%%%%%%%%%%%%%%%%%%%%%%%%%%%%%%%%%%%%%%%%%%%%%%%%%%%  УПРАВЛЕНИЕ ВЫКАТНЫМ ЭЛЕМЕНТОМ %%%%%%%%%%%%%%%%%%%%%%%%%%%%%%%%%%%%%%%%%%%%%%%%%%

\item Управление выкатным элементом (УВЭ)

\begin{enumerate}[label=\arabic{section}.\arabic{subsection}.\arabic{enumi}.\arabic*, labelsep=4pt, leftmargin=0em, itemindent=65pt, parsep=0pt]

\item 
В функциональном блоке «КП» логика управления выкатным элементом представлена функцией «УВЭ». Функция формирует команды вкатить и выкатить выкатной элемент ячейки фидера при оперативном управлении.

\item
Команда операции выкатить ВЭ «УВЭ: Выкатить» формируется в следующих случаях:
\begin{itemize}
	\item при появлении команды выкатить от местного управления («Выкат. ВЭ ИЧМ»), выбранным управлением по месту (сигнал «Местное» находится в активном состоянии) и наличием разрешающего сигнала от схемы оперативной блокировки «ОБР ВЭ: Разрешить управ.»;
	\item при появлении команды выкатить от дистанционного управления («Выкат. ВЭ АСУ»), выбранным дистанционным управлением (сигнал «Дистанционное» находится в активном состоянии) и наличием разрешающего сигнала от схемы оперативной блокировки «ОБР ВЭ: Разрешить управ.».
\end{itemize}

\item
Прохождение команды выкатить ВЭ блокируется при наличии следующих сигналов:
\begin{itemize}
	\item контрольное положение ВЭ;
	\item превышение времени операции вкатить/выкатить ВЭ «УВЭ: Прев. времени~пер.».
\end{itemize}

\item
Команда операции вкатить ВЭ «УВЭ: Вкатить» формируется в следующих случаях:
\begin{itemize}
	\item при появлении команды вкатить от местного управления («Вкат. ВЭ ИЧМ»), выбранным управлением по месту (сигнал «Местное» находится в активном состоянии) и наличием разрешающего сигнала от схемы оперативной блокировки «ОБР ВЭ: Разрешить управ.»;
	\item при появлении команды вкатить от дистанционного управления («Вкат. ВЭ АСУ»), выбранным дистанционным управлением (сигнал «Дистанционное» находится в активном состоянии) и наличием разрешающего сигнала от схемы оперативной блокировки «ОБР ВЭ: Разрешить управ.».
\end{itemize}

\item
Прохождение команды вкатить ВЭ блокируется при наличии следующих сигналов:
\begin{itemize}
	\item рабочее положение ВЭ;
	\item превышение времени операции вкатить/выкатить ВЭ «УВЭ: Прев. времени~пер.».
\end{itemize}

\item
При наличии сигналов «УВЭ: Вкатить» и «УВЭ: Выкатить» длительностью больше, чем время уставки «T1» формируется сигнал превышения допустимого времени операции -- «УВЭ: Прев. времени пер.». Активное состояние операции включения или отключения сигнализируется выходным сигналом «УВЭ: Идет переключение».

\item
Логика функции формирует сигналы состояния выключателя «УВЭ: Выкачен», «УВЭ: Вкачен», «УВЭ: Не определено» и «УВЭ: Неисправность», при этом сигналы «УВЭ: Неисправность» и «УВЭ: Не определено» выдаются с выдержкой времени, задаваемой уставкой Т2 «Тблк».

\item
Ввод функции в работу на этапе параметрирования устройства осуществляется программным переключателем SGF1 «Ввод\_функции». Информация о введенной функции отображается выходным сигналом «УВЭ: Ввод».

\item
Алгоритм работы выполнен в соответствии с рисунком \ref{kp:img_uve}. 

\vspace{1mm}
\begin{figure}[H]
\centering
\includegraphics[width=1\textwidth,height=1\textheight,keepaspectratio]{\fbpath/128. КП ОЛ/_latex/img/uve.pdf}
\captionof{figure}{Функциональная схема логики УВЭ}\label{kp:img_uve}
\end{figure}


\end{enumerate}
%%%%%%%%%%%%%%%%%%%%%%%%%%%%%%%%%%%%%%%%%%%%%%%%%%%%%%%%%%% КОНЕЦ УПРАВЛЕНИЕ ВЫКАТНЫМ ЭЛЕМЕНТОМ %%%%%%%%%%%%%%%%%%%%%%%%%%%%%%%%%%%%%%%%%%%%%%%%%%

%%%%%%%%%%%%%%%%%%%%%%%%%%%%%%%%%%%%%%%%%%%%%%%%%%%%%%%%%%%%%  УПРАВЛЕНИЕ ЗАЗЕМЛЯЮЩИМ НОЖОМ %%%%%%%%%%%%%%%%%%%%%%%%%%%%%%%%%%%%%%%%%%%%%%%%%%
\item Управление заземляющим ножом (УЗН)

\begin{enumerate}[label=\arabic{section}.\arabic{subsection}.\arabic{enumi}.\arabic*, labelsep=4pt, leftmargin=0em, itemindent=65pt, parsep=0pt]

\item 
В функциональном блоке «КП» логика управления заземляющим ножом представлена функцией «УЗН». Функция формирует команды на отключение и включение заземляющего ножа при оперативном управлении.

\item
Команда оперативного отключения «УЗН: Отключить» формируется в следующих случаях:
\begin{itemize}
	\item при появлении команды на оперативное отключение заземляющего ножа («Откл. ЗН ИЧМ»), выбранным управлением по месту (сигнал «Местное» находится в активном состоянии) и наличием разрешающего сигнала от схемы оперативной блокировки «ОБР ЗН: Разрешить управ.»;
	\item при появлении команды на оперативное отключение заземляющего ножа («Откл. ЗН АСУ»), выбранным дистанционным управлением (сигнал «Дистанционное» находится в активном состоянии) и наличием разрешающего сигнала от схемы оперативной блокировки «ОБР ЗН: Разрешить управ.».
\end{itemize}

\item
Прохождение команды на отключение блокируется при наличии следующих сигналов:
\begin{itemize}
	\item отключенное положение заземляющего ножа «ЗН отключен (б/к)» (с контролем сигнала «ЗН включен~(б/к)»);
	\item превышение времени переключения заземляющего ножа «УЗН: Прев. времени~пер.».
\end{itemize}

\item
Команда оперативного включения «УЗН: Включить» формируется в следующих случаях:
\begin{itemize}
	\item при появлении команды на оперативное включение выключателя («Вкл. ЗН ИЧМ»), выбранным управлением по месту (сигнал «Местное» находится в активном состоянии) и наличием разрешающего сигнала от схемы оперативной блокировки «ОБР ЗН: Разрешить управ.»;
	\item при появлении команды на оперативное включение выключателя («Вкл. ЗН АСУ»), выбранным дистанционным управлением (сигнал «Дистанционное» находится в активном состоянии) и наличием разрешающего сигнала от схемы оперативной блокировки «ОБР ЗН: Разрешить управ.».
\end{itemize}

\item
Прохождение команды на включение блокируется при наличии следующих сигналов:
\begin{itemize}
	\item включенное положение заземляющего ножа «ЗН~включен (б/к)» (с контролем сигнала «ЗН~отключен~(б/к)»);
	\item превышение времени переключения заземляющего ножа «УЗН: Прев. времени~пер.».
\end{itemize}

\item
При наличии сигналов «УЗН: Включить» и «УЗН: Отключить» длительностью больше, чем время уставки «T1» формируется сигнал превышения допустимого времени переключения -- «УЗН: Прев. времени пер.». Активное состояние операции включения или отключения сигнализируется выходным сигналом «УЗН: Идет переключение».

\item
Логика функции формирует сигналы состояния выключателя «УЗН: Отключено», «УЗН: Включено», «УЗН: Не определено» и «УЗН: Неисправность», при этом сигналы «УЗН: Неисправность» и «УЗН: Не определено» выдаются с выдержкой времени, задаваемой уставкой Т2 «Тблк».

\item
Ввод функции в работу на этапе параметрирования устройства осуществляется программным переключателем SGF1 «Ввод\_функции». Информация о введенной функции отображается выходным сигналом «УЗН: Ввод».

\item
Алгоритм работы выполнен в соответствии с рисунком \ref{kp:img_uzn}. 

\vspace{1mm}
\begin{figure}[H]
\centering
\includegraphics[width=1\textwidth,height=1\textheight,keepaspectratio]{\fbpath/128. КП ОЛ/_latex/img/uzn.pdf}
\captionof{figure}{Функциональная схема логики УЗН}\label{kp:img_uzn}
\end{figure}


\end{enumerate}

%%%%%%%%%%%%%%%%%%%%%%%%%%%%%%%%%%%%%%%%%%%%%%%%%%%%%%%%%%%% КОНЕЦ УПРАВЛЕНИЕ ЗАЗЕМЛЯЮЩИМ НОЖОМ %%%%%%%%%%%%%%%%%%%%%%%%%%%%%%%%%%%%%%%%%%%%%%%%%%

%%%%%%%%%%%%%%%%%%%%%%%%%%%%%%%%%%%%%%%%%%%%%%%%%%%%%%%%%%%%  ОПЕРАТИВНАЯ БЛОКИРОВКА ОБЩЕЕ %%%%%%%%%%%%%%%%%%%%%%%%%%%%%%%%%%%%%%%%%%%%%%%%%%
\item Оперативная блокировка КА (ОБР КА)

\begin{enumerate}[label=\arabic{section}.\arabic{subsection}.\arabic{enumi}.\arabic*, labelsep=4pt, leftmargin=0em, itemindent=65pt, parsep=0pt]

\item
В устройстве реализованы алгоритмы, которые позволяют контролировать положение КА и выдавать сигналы на разрешение управления выключателем (В), выкатным элементом выключателя (ВЭ) и заземляющим ножом (ЗН).

\item
Поясняющая схема первичной сети для алгоритмов оперативной блокировки КА приведена на рисунке \ref{kp:img_obrka}.

\vspace{1mm}
\begin{figure}[H]
\centering
\includegraphics[width=0.5\textwidth,height=0.5\textheight,keepaspectratio]{\fbpath/128. КП ОЛ/_latex/img/obrka.jpg}
\captionof{figure}{Поясняющая схема первичной сети для алгоритмов оперативной блокировки КА}\label{kp:img_obrka}
\end{figure}

\end{enumerate}
%%%%%%%%%%%%%%%%%%%%%%%%%%%%%%%%%%%%%%%%%%%%%%%%%%%%%%%%%%%%%%%%%%%%%% КОНЕЦ ОПЕРАТИВНАЯ БЛОКИРОВКА ОБЩЕЕ %%%%%%%%%%%%%%%%%%%%%%%%%%%%%%%%%%%%%%%%%%%%%%%%%%

%%%%%%%%%%%%%%%%%%%%%%%%%%%%%%%%%%%%%%%%%%%%%%%%%%%%%%%%%%%%%%%%%%%%%%  Оперативная блокировка выключателя (ОБР В) %%%%%%%%%%%%%%%%%%%%%%%%%%%%%%%%%%%%%%%%%%%%%%%%%%
\item Оперативная блокировка выключателя (ОБР В)

\begin{enumerate}[label=\arabic{section}.\arabic{subsection}.\arabic{enumi}.\arabic*, labelsep=4pt, leftmargin=0em, itemindent=65pt, parsep=0pt]

\item 
В функциональном блоке «КП» логика оперативной блокировки выключателя представлена функцией «ОБР В». Условием ввода в работу разрешения управления выключателем является наличие входного сигнала «УВ: Ввод». Контроль условий блокировки управления выключателем осуществляется с помощью уставки SGF1. Если уставка SGF1 выставлена в положение «Не предусмотрено», то управление выключателем осуществляется без учета всех блокировок.

\item
Если уставка SGF1 выставлена в положение «Предусмотрено», то разрешение управления выключателем формируется при выполнении следующих условий:
\begin{itemize}
	\item отсутствие операции переключения коммутационных аппаратов (сигнал «КП: Идет переключение»);
	\item активное состояние сигнала «ВЭ: Выкачен» для управления выключателем в контрольном (ремонтном) положениях или одновременное наличие следующих сигналов для разрешения 	управления выключателем в рабочем положении;
	\item активное состояние сигналов: «ВЭ: Вкачен», «ЗН: Отключен», «КА1 Отключен», «КА2 Отключен», если введен контроль данных сигналов соответствующими программными переключателями.	
\end{itemize}

\item\sloppy
Внешние дискретные входы назначаются в проекте непосредственно под схему объекта. Блокировка включения выключателя выполняется при несоблюдении какого-либо из условий, предусмотренных алгоритмом функции.

\item
Алгоритм работы функции выполнен в соответствии с рисунком \ref{kp:img_obrv}. 

\vspace{1mm}
\begin{figure}[H]
\centering
\includegraphics[width=0.85\textwidth,height=0.85\textheight,keepaspectratio]{\fbpath/128. КП ОЛ/_latex/img/obrv.pdf}
\captionof{figure}{Функциональная схема логики ОБР В}\label{kp:img_obrv}
\end{figure}


\end{enumerate}
%%%%%%%%%%%%%%%%%%%%%%%%%%%%%%%%%%%%%%%%%%%%%%%%%%%%%%%%%%%%%%%%%%%%%% КОНЕЦ Оперативная блокировка выключателя (ОБР В) %%%%%%%%%%%%%%%%%%%%%%%%%%%%%%%%%%%%%%%%%%%%%%%%%%

%%%%%%%%%%%%%%%%%%%%%%%%%%%%%%%%%%%%%%%%%%%%%%%%%%%%%%%%%%%%%%%%%%%%%%  Оперативная блокировка выкатного элемента (ОБР ВЭ) %%%%%%%%%%%%%%%%%%%%%%%%%%%%%%%%%%%%%%%%%%%%%%%%%%
\item Оперативная блокировка выкатного элемента (ОБР ВЭ)

\begin{enumerate}[label=\arabic{section}.\arabic{subsection}.\arabic{enumi}.\arabic*, labelsep=4pt, leftmargin=0em, itemindent=65pt, parsep=0pt]

\item 
В функциональном блоке «КП» логика оперативной блокировки выкатного элемента представлена функцией «ОБР ВЭ». Условие ввода в работу разрешения управления ВЭ является наличие входного сигнала «УВЭ: Ввод». Контроль условий блокировки управления ВЭ осуществляется с помощью уставки SG1. Если уставка SG1 (см. таблицу \ref{kp:tbl_gen}) выставлена в положение «Не предусмотрено», то управление ВЭ осуществляется без учета всех блокировок.

\item
Принудительное деблокирование управления выкатным элементом от оперативной блокировки осуществляется активацией входного сигнала от ключа «Программ. деблок. ВЭ».

\item
Разрешение управления ВЭ формируется при выполнении необходимых условий:
\begin{itemize}
	\item отсутствие процесса переключения коммутационных аппаратов (сигнал «КП: Идет переключение»);
	\item активное состояние сигнала «ВЭ: Вкачен» для разрешения команды выкатить ВЭ или одновременное наличие следующих сигналов для разрешения команды вкатить ВЭ: «ВЭ: Выкачен», «ЗН: Отключен», «КА1 Отключен», «КА2 Отключен», если введен контроль данных сигналов соответствующими уставками.
\end{itemize}

\item\sloppy
Внешние дискретные входы назначаются в проекте непосредственно под схему объекта. Блокировка управления ВЭ выполняется при несоблюдении какого-либо из условий, предусмотренных алгоритмом функции.

\item
Алгоритм работы функции выполнен в соответствии с рисунком \ref{kp:img_obrve}. 

\vspace{1mm}
\begin{figure}[H]
\centering

\includegraphics[width=1\textwidth,height=1\textheight,keepaspectratio]{\fbpath/128. КП ОЛ/_latex/img/obrve.pdf}
\captionof{figure}{Функциональная схема логики ОБР ВЭ}\label{kp:img_obrve}
\end{figure}

\end{enumerate}
%%%%%%%%%%%%%%%%%%%%%%%%%%%%%%%%%%%%%%%%%%%%%%%%%%%%%%%%%%%%%%%%%%%%%% КОНЕЦ Оперативная блокировка выкатного элемента (ОБР ВЭ) %%%%%%%%%%%%%%%%%%%%%%%%%%%%%%%%%%%%%%%%%%%%%%%%%%

%%%%%%%%%%%%%%%%%%%%%%%%%%%%%%%%%%%%%%%%%%%%%%%%%%%%%%%%%%%%%%%%%%%%%%  Оперативная блокировка заземляющего ножа (ОБР ЗН) %%%%%%%%%%%%%%%%%%%%%%%%%%%%%%%%%%%%%%%%%%%%%%%%%%
\item Оперативная блокировка заземляющего ножа (ОБР ЗН)

\begin{enumerate}[label=\arabic{section}.\arabic{subsection}.\arabic{enumi}.\arabic*, labelsep=4pt, leftmargin=0em, itemindent=65pt, parsep=0pt]

\item
В функциональном блоке «КП» логика оперативной блокировки заземляющего ножа представлена функцией «ОБР ЗН». Условие ввода в работу разрешения управления ЗН является наличие входного сигнала «УЗН: Ввод». Контроль условий блокировки управления ВЭ осуществляется с помощью уставки SG1. Если уставка SG1 (см. таблицу \ref{kp:tbl_gen}) выставлена в положение «Не предусмотрено», то управление ЗН осуществляется без учета всех блокировок.

\item
Принудительное деблокирование управления заземляющим ножом от оперативной блокировки осуществляется активацией входного сигнала от ключа «Программ. деблок. ЗН».

\item
Разрешение управления ЗН формируется при выполнении необходимых условий:
\begin{itemize}
	\item отсутствие процесса переключения коммутационных аппаратов (сигнал «КП: Идет переключение»);
	\item активное состояние сигнала «ЗН: Включен» для разрешения команды отключения ЗН или одновременное наличие следующих сигналов для разрешения команды включения ЗН: «ВЭ: Выкачен», «ЗН: Отключен», «КА3 Отключен», «КА4 Отключен», если введен контроль данных сигналов соответствующими уставками.
\end{itemize}

\item\sloppy
Внешние дискретные входы назначаются в проекте непосредственно под схему объекта. Блокировка управления ЗН выполняется при несоблюдении какого-либо из условий, предусмотренных алгоритмом функции.

\item
Алгоритм работы функции выполнен в соответствии с рисунком \ref{kp:img_obrzn}. 

\vspace{1mm}
\begin{figure}[H]
\centering

\includegraphics[width=1\textwidth,height=1\textheight,keepaspectratio]{\fbpath/128. КП ОЛ/_latex/img/obrzn.pdf}
\captionof{figure}{Функциональная схема логики ОБР ЗН}\label{kp:img_obrzn}
\end{figure}


\end{enumerate}
%%%%%%%%%%%%%%%%%%%%%%%%%%%%%%%%%%%%%%%%%%%%%%%%%%%%%%%%%%%%%%%%%%%%%% КОНЕЦ Оперативная блокировка заземляющего ножа (ОБР ЗН) %%%%%%%%%%%%%%%%%%%%%%%%%%%%%%%%%%%%%%%%%%%%%%%%%%

\end{enumerate}
